\subsubsection{Tiền xử lý}

\paragraph{Xử lý các giá trị còn thiếu}

Bộ dữ liệu Forest CoverType không có giá trị thiếu nên bước này được bỏ qua.

\paragraph{Chuẩn hóa dữ liệu} Các biến đặc trưng được chuẩn hoá về trung bình
$\mu = 0$ và phương sai $\sigma^2 = 1$. Điều này giúp các giá trị được đưa về
cùng thang đo, giúp quá trình học bằng Gradient ổn định và nhanh hơn.

\paragraph{Chia dữ liệu}

Dữ liệu được chia thành 3 tập là huấn luyện, xác thực và tập kiểm tra với tỉ lệ
tương ứng là 0.6, 0.2 và 0.2. Do dữ liệu bị mất cân bằng nên quá trình chia đã
được chỉnh để giữ tỉ lệ giá trị các lớp của biến mục tiêu trên cả 3 tập.
